\subsection{Расширение на сложные структуры и независимая оценка}
Расширение сохраняет задачу подготовки реакции, варьируя управление потоком и взаимодействие ресурсов (табл.~\ref{tab:processes}). P1 синхронизирует две операции с общим пулом. P2 направляет простые и сложные экземпляры по разным маршрутам и сохраняет выбранную ветвь при возврате. P3 сочетает локальный возврат с повторением всего процесса. P4 содержит отдельные потоки поступлений, средние операции которых конкурируют за общий ресурс. В P5 два узких места могут по-разному реагировать на добавление мощности. P6 объединяет маршрутизацию по классу, параллельное исполнение и повтор выбранной ветви.

\begin{table}[t]\caption{Семейства процессов. Указаны числа наблюдаемых операций, ресурсных пулов, отдельных точек возврата и типов вмешательств; отсутствие действия доступно для каждого семейства.}\label{tab:processes}\centering\small\setlength{\tabcolsep}{4pt}
\begin{tabular}{lp{.40\textwidth}rrrr}\toprule
 & Структура & Опер. & Пулы & Циклы & Действия\\\midrule
P0 & Последовательность с возвратом & 3 & 3 & 1 & 2\\
P1 & AND-блок, синхронизация, повтор & 4 & 3 & 1 & 2\\
P2 & XOR по классу, повтор ветви & 5 & 4 & 2 & 3\\
P3 & Локальный и глобальный циклы & 4 & 4 & 2 & 3\\
P4 & Два процесса, общий средний пул & 6 & 5 & 2 & 3\\
P5 & Узкие места, возврат к поздней стадии & 5 & 5 & 1 & 3\\
P6 & XOR с вложенным AND, повтор ветви & 7 & 6 & 2 & 3\\
\bottomrule\end{tabular}\end{table}

Каждая новая структура оценивается в мирах fixed, service, routing и combined: по 12 независимых обучающих наборов на сочетание, всего 288. Истинное замедление обслуживания задаётся насыщающейся функцией длины очереди; вероятность возврата — насыщающейся квадратичной функцией глобальной, локальной или процессной нагрузки. Генераторы, распределение ресурсов, действия и стоимости фиксируются до сравнения методов. Калибровка нормальной нагрузки использует только генератор и заранее заданную сетку проверки устойчивости, без учёта результатов методов. Архив содержит полные спецификации и все просмотренные калибровочные исходы.

Пороги нагрузки масштабируются относительно нормального режима. $S$ — округлённый вверх 95-й процентиль WIP в отдельных опорных эпизодах мира fixed; $T$ — округлённая вверх средняя длительность экземпляра. Для P1--P6 получены $S=(3,4,5,5,5,5)$ и $T=(7,7,8,7,10,8)$. Кандидатные пороги равны $2S,4S,8S$, интервал наблюдений — $0.5T$, тяжёлая перегрузка требует WIP не ниже $10S$ в течение $20T$. Отдельная сетка определяет устойчивую опорную интенсивность для каждого процесса. Обучение содержит восемь эпизодов при $0.4,0.6,0.8,1.0$ от этой интенсивности, по два на уровень; валидация использует $0.5,0.7,0.9$. Поступления завершаются в $40T$, после чего система разгружается. Все обучающие и валидационные эпизоды проверяются по критерию тяжёлого события. P0 сохраняет абсолютные пороги и исходный протокол.

Планирование использует шесть условий с двумя повторами и горизонтом $140T$; тестирование — шесть других условий с восемью повторами и горизонтом $180T$. В тестовых импульсах меняются время и интенсивность. Калибровка, валидация, планирование и тест используют раздельные случайные потоки. Все кандидатные действия проверяются на истинном процессе после фиксации политик. Варианты с известной структурой и discovery только по обучению используют одинаковые журналы, каталог действий и парные тестовые траектории. Fitness, precision, добавленные и пропущенные пути, отказы исполнимой модели регистрируются без исправления обнаруженной структуры.

Помимо static, predictive selection и retained-load, расширение проверяет абляции только времени обслуживания и только маршрутизации, отключение отдельных точек возврата, каждую фиксированную пару действия и порога, а также планировщик с известным генератором. Эвристика силы обратной связи выбирает действие по оценённому локальному усилению длительности или повторений и использует порог $2S$; её входные данные получены из обучающих журналов, без имитационного сравнения альтернатив. Regret относительно планировщика с известным генератором сравнивает реализованные стоимости политик, выбранных по одной конечной библиотеке планирования. Минимум по тесту приводится только как ретроспективная граница.

Основные новые сравнения сопоставляют retained-load с predictive selection при routing feedback отдельно для P1--P6 с известной структурой. Траектории усредняются внутри обученной модели, затем используются парные различия по 12 независимым обучениям. Поправка Бонферрони охватывает шесть основных сравнений; остальные интервалы описательные. При межпроцессном объединении каждый процесс получает равный вес после нормировки на стоимость no action. Известные и обнаруженные структуры рассматриваются раздельно. Диагностические тесты варьируют пороги селектора, меняют долю сложных экземпляров P2 с $0.35$ на $0.65$ и фактическое сокращение rework в P3 на $30\%$ или $70\%$, сохраняя политики, выбранные при номинальном эффекте. Эти тесты не перенастраивают метод.
